\section*{Chapter \ref{ch:qm:volatility-models} --- Volatility Models}

\begin{solution}{exo:qm:volatility-models:1}
$\bar h = 0.02/(1 - 0.98) = 1$: daily volatility 1\%, annual $\sqrt{252} = 15.9\%$; \omterm{def:qm:stochastic-differential-equations:ou}{half-life} $\ln\frac12/\ln0.98 = 34.3$ days.
\end{solution}

\begin{solution}{exo:qm:volatility-models:2}
$\ln\frac12/\ln0.94 = 11.2$ days. The last 20 days carry $1 - 0.94^{20} = 71\%$ of the weight.
\end{solution}

\begin{solution}{exo:qm:volatility-models:3}
$m = 24 \times 12 = 288$ and $\sqrt{2/288} = 8.3\%$.
\end{solution}

\begin{solution}{exo:qm:volatility-models:4}
$h_{t+1} = 4\bar h$ and $\E_t[h_{t+k}] = \bar h + 0.98^{k-1} \times 3\bar h$; averaging over $k = 1, \dots, 10$ gives $3.74\bar h$, a volatility of $\sqrt{3.74} = 1.93\%$
a day. The forecast barely reverts in ten days when the \omterm{def:qm:stochastic-differential-equations:ou}{half-life} is 34.
\end{solution}

\begin{solution}{exo:qm:volatility-models:5}
$\E[\hat\sigma^2/h + \ln h] = \sigma^2/h + \ln h$, whose derivative $-\sigma^2/h^2 + 1/h$ vanishes at $h = \sigma^2$, a minimum. With $\sigma^2 = 1$, the loss is 1 at the
truth, $2 + \ln0.5 = 1.31$ for $0.5\sigma^2$ and $0.67 + \ln1.5 = 1.07$ for $1.5\sigma^2$: under-prediction is punished more.
\end{solution}

\begin{solution}{exo:qm:volatility-models:6}
With the 3.5\% return, the variance (about zero mean) is $(249 \times 0.25 + 12.25)/250 = 0.298$, a volatility of 0.546\%; without it, 0.5\%. The window's
volatility falls by 8.4\% on the day the return leaves, whatever happens that day.
\end{solution}

\begin{solution}{exo:qm:volatility-models:7}
With the normal likelihood: $\alpha = 0.029$ with Hessian \omterm{def:qm:estimation:estimator}{standard error} 0.0026 and sandwich 0.0033; $\beta = 0.968$ with 0.0028 and 0.0035. The sandwich is about 25\%
larger because the standardised residuals have fat tails, so the normal likelihood is misspecified and the information equality fails; report the sandwich.
With the Student-$t$ likelihood the two agree (0.0033 and 0.0032 for $\alpha$).
\end{solution}

\begin{solution}{exo:qm:volatility-models:8}
The absolute return is a biased proxy for volatility ($\E|r| = \sqrt{2/\pi}\,\sigma$ for normal returns, less for fat tails), and mean absolute error is not in Patton's
robust family: under a noisy proxy it can rank a biased forecast above the true variance, favouring forecasts that are too low. Compare variance forecasts
with squared error or QLIKE against squared returns or \omterm{def:qm:volatility-models:rv}{realised variance}, and test the difference.
\end{solution}

\begin{solution}{pb:qm:volatility-models:1}
\textbf{1.} $\omega = 0.0010$, $\alpha = 0.029$ (0.003), $\beta = 0.968$ (0.004).
\textbf{2.} $\omega = 0.0007$, $\alpha = 0.030$ (0.003), $\beta = 0.968$ (0.003), $\nu = 7.2$ (0.6); log-likelihood higher by 141.
\textbf{3.} Persistence 0.9984; \omterm{def:qm:stochastic-differential-equations:ou}{half-life} 445 trading days (255 with the normal likelihood).
\textbf{4.} $\sqrt{\omega/(1 - \alpha - \beta)} = 0.67\%$ against a sample standard deviation of 0.58\%: with persistence this close to one, the implied level is poorly
determined.
\textbf{5.} No: the GJR asymmetry is 0.006 with \omterm{def:qm:estimation:estimator}{standard error} 0.005.
\textbf{6.} The 3.5\% rise of 11 November 2022, leaving the window at the close of 2 November 2023.
\textbf{7.} From 0.538\% to 0.496\%, by 7.8\%, while the euro rose 1.2\% that day.
\textbf{8.} GARCH rose from 0.424\% to 0.465\%, EWMA from 0.410\% to 0.490\%.
\textbf{9.} Under GARCH its effect halves in 445 days; under EWMA in 11 days; in the window it stays whole for 250 days and vanishes at once.
\textbf{10.} That risk fell, on a day of an unusually large move.
\textbf{11.} QLIKE $-0.521$ (GARCH-$t$ fitted on 1999--2014), $-0.495$ (EWMA), $-0.435$ (window), over 3\,002 days.
\textbf{12.} Yes: Diebold--Mariano $-2.76$ (GARCH against EWMA, $p = 0.006$), $-4.10$ (GARCH against the window), $-2.48$ (EWMA against the window).
\textbf{13.} 0.77, 0.60 and 0.56: all below one, so every forecast moves too much with recent data; GARCH least.
\textbf{14.} The proxy, a squared daily return, is the variance times a $\chi^2_1$-like variable, so most of its variation is noise no forecast can explain; $R^2$ of 3\% is
near the ceiling.
\textbf{15.} A far less noisy proxy: sharper rankings, higher $R^2$, and more power in the \omterm{def:qm:volatility-models:eval}{Diebold--Mariano tests}.
\textbf{16.} The GARCH-$t$ forecast, or EWMA if simplicity matters more than the 5\% gain in QLIKE; never the equally weighted window.
\textbf{17.} As a description of how slowly large shocks fade, yes; as a precise number, no: persistence near one is the short-memory model's way of approximating
\omterm{def:qm:linear-time-series:longmem}{long memory}, and the \omterm{def:qm:stochastic-differential-equations:ou}{half-life}'s confidence interval is wide.
\textbf{18.} That the fall in reported risk was an artefact of the window's calendar, and that every forecast that weights by time rose.
\textbf{19.} Named result: \emph{the ghost in the risk report}: a EUR/USD variance shock has a \omterm{def:qm:stochastic-differential-equations:ou}{half-life} of 445 trading days under GARCH(1,1) with Student-$t$
innovations; the 250-day window's volatility fell 7.8\% on 2 November 2023, the day the 3.5\% return of 11 November 2022 left it.
\textbf{20.} It keeps every return at full weight for a fixed time and then drops it, so the forecast moves on anniversaries rather than on news.
\end{solution}

\begin{solution}{iq:qm:volatility-models:1}
Because a large return leaves the window that day: the fall is the anniversary of old news, not information about today, and it can outweigh a large new return.
\iqlookfor{The drop-off effect and its fix (decaying weights).}
\end{solution}

\begin{solution}{iq:qm:volatility-models:2}
$h_t = \omega + \alpha r_{t-1}^2 + \beta h_{t-1}$; unconditional variance $\omega/(1 - \alpha - \beta)$ when $\alpha + \beta < 1$; $\alpha + \beta$ is the persistence, the daily decay rate
of a variance shock, with \omterm{def:qm:stochastic-differential-equations:ou}{half-life} $\ln\frac12/\ln(\alpha + \beta)$. $\alpha$ is the reaction to news, $\beta$ the memory.
\iqlookfor{The formula, stationarity, and the persistence as a decay rate.}
\end{solution}

\begin{solution}{iq:qm:volatility-models:3}
EWMA is GARCH(1,1) with $\omega = 0$ and $\alpha + \beta = 1$: integrated. It forecasts today's variance at every horizon, while GARCH reverts to its long-run level; for
long horizons after a shock, EWMA overstates (or, after a quiet spell, understates) the variance.
\iqlookfor{IGARCH and the flat term structure of forecasts.}
\end{solution}

\begin{solution}{iq:qm:volatility-models:4}
Compare against an unbiased proxy (squared returns or \omterm{def:qm:volatility-models:rv}{realised variance}) with a loss that ranks correctly under proxy noise, squared error or QLIKE, test the loss
difference with Diebold--Mariano, and check calibration with a \omterm{def:qm:volatility-models:eval}{Mincer--Zarnowitz regression}.
\iqlookfor{Proxy, robust loss, and a test, not a single number.}
\end{solution}

\begin{solution}{iq:qm:volatility-models:5}
The sum of squared intraday returns, which converges to the \omterm{def:qm:volatility-models:rv}{integrated variance} as the sampling gets finer. At very high frequency, bid-ask bounce and price
discreteness add noise whose contribution grows with the number of returns, so RV explodes; sample at a few minutes, or use noise-robust \omterm{def:qm:estimation:estimator}{estimators}
(chapter~21).
\iqlookfor{Convergence, and the noise trade-off.}
\end{solution}

\begin{solution}{iq:qm:volatility-models:6}
Its three regressors average \omterm{def:qm:volatility-models:rv}{realised variance} over a day, a week and a month, so the forecast is a sum of components with short, medium and \omterm{def:qm:linear-time-series:longmem}{long memory}; that
cascade mimics the slowly decaying, long-memory autocorrelation of volatility, while staying a linear regression that is easy to fit and hard to overfit.
\iqlookfor{The multiple horizons as an approximation to \omterm{def:qm:linear-time-series:longmem}{long memory}.}
\end{solution}
